\section{Discussion}

\subsection{Formulation Choices and Conservative Structure}

The central modelling choice in this work is to use one fixed porosity field,
$\theta$, in the finite-volume transport equation and in the PHREEQC cell
volumes alike, which keeps aqueous accumulation, transport, and equilibrium
chemistry on the same volumetric basis. The reactive surface is prescribed
independently through the explicit cell field $a_{\mathrm{surf}}$, so the
chemical storage volume and the amount of reactive surface remain separately
defined material properties.

This approach extends earlier micro-continuum reactive transport formulations
to the case where the reactive surface area is an unresolved cell property
rather than an explicitly meshed boundary. In existing pore-scale DBS solvers,
including \nolinkurl{porousMedia4Foam} and earlier GeoChemFoam solvers,
reactive uptake is typically imposed as a flux condition on the diffuse
fluid-solid interface. This requires the reactive surface to coincide with the resolved image voxel
boundary, which is a good approximation when the dominant reactive area lies at
the contact between macropore and grain. In multi-scale carbonate images the
microporous phase holds a large internal reactive area that lies inside the
coarse cell rather than on its boundary. The cell-based formulation carries
that area as an intrinsic volumetric density
$a_{\mathrm{surf}}^{\mathrm{intrinsic}}$, scaled by the local phase fraction
and passed to PHREEQC as a cell-level property.

The closed-batch balance isolates the round-trip transfer of aqueous and
adsorbed Ca between the finite-volume fields and PhreeqcRM: with transport and
boundary fluxes disabled, the combined inventory
$M=\theta C+a_{\mathrm{surf}}\Gamma$ should remain unchanged during local
equilibration. The transient 1D benchmark then tests the coupled transport
response, while its mesh- and timestep-refinement study quantifies the combined
spatial, temporal, and sequential-coupling errors.

\subsection{Geometric Scalability}

Because the validation ladder separates aqueous--surface state transfer,
coupled transport, surface-area geometry, and microporosity, the evidence can
be attributed to the layer that produced it. The closed-batch residual remains
below $10^{-9}$ in relative magnitude, the reported transient profiles differ
from the finest combined result by $3.01\%$ and $1.05\%$ for the
surface-enabled and aqueous-only cases, and the interface-area construction is
accurate to sub-percent levels against analytic spheres at
$R/\Delta x\ge16$ and against circular posts at all four tested resolutions. The
intrinsic micropore density enters only in the cases that use it.

The transition from 2D to 3D introduces qualitatively different transport paths
and surface-loading distributions. In the 2D case, open fractures provide
preferential pathways for dissolved Pb$^{2+}$, while strong-site Pb occupancy
varies across the surrounding microporous matrix and is enhanced near the inlet
and fracture intersections. In the 3D cases, the tortuosity of the
pore network produces a more dispersed front and a more heterogeneous
surface-loading pattern. This spatial heterogeneity in surface uptake controls the effective retardation
factor seen at the continuum scale, and following it requires a simulation that
resolves the local connectivity of the image.

The direct resolved-to-DBS comparison spans $Pe=0.1$ to $10$, with no
observed breakdown in agreement at the higher $Pe$ values tested. This
behaviour is consistent with the spatially uniform matrix closure and
repeated geometry of the benchmark, whose limited heterogeneity may help
preserve the accuracy of the two-scale representation as advection becomes
more dominant.

\subsection{Computational Cost Drivers}

PHREEQC is called once per active cell per timestep, with a full speciation
and surface-equilibration calculation for each cell. The chemistry overhead
therefore depends primarily on the number of active cells, the geochemical
system, and the frequency of the chemistry call. The resolved-wall and DBS
comparisons use the same nominal image resolution and quantify physical
agreement. A controlled performance comparison across resolutions is outside
the scope of this study. For larger grids or more complex geochemical systems,
masking near-solid cells
out of the PhreeqcRM mapping (here at
$\theta \le \theta_{\mathrm{mask}} = 10^{-3}$) is the main lever on that cost.
The threshold is held fixed across all examples here.

\subsection{Limitations and Future Work}

The intrinsic micropore end member $s_{\mathrm{micro}}$ is prescribed from
complementary measurements such as nitrogen BET adsorption, mercury intrusion
porosimetry, or finer-resolution image analysis \citep{noiriel2009changes}.
The segmented image determines the geometric fraction $f_{\mathrm{micro}}$,
and their product gives the bulk-volume cell contribution
$a_{\mathrm{surf}}^{\mathrm{intrinsic}}=f_{\mathrm{micro}}s_{\mathrm{micro}}$.
Uncertainty in $s_{\mathrm{micro}}$ is therefore a leading physical uncertainty
in microporous applications. This uncertainty is not specific to the
cell-based formulation: a resolved-wall representation constructed from the
same segmented image also cannot recover reactive area below the image
resolution without complementary characterisation. Linking image-derived
phase fractions to measured or computed phase-specific surface densities is
therefore required to constrain predicted uptake.

The upscaled DBS treatment carries unresolved matrix properties as local
cell-averaged quantities. When diffusive exchange is fast compared with the
advective residence time, the intrinsic surface area acts as a locally
accessible reactive reservoir. At larger $Pe$, advective bypass and finite
matrix-exchange times can restrict access to that area, so reactive uptake
depends on its time-dependent accessibility as well as its amount. This
transport-upscaling limitation is distinct from the assumption of local
chemical equilibrium within each active cell: local equilibration does not
ensure that solute reaches the unresolved reactive surface. Media with more
heterogeneous matrix properties, connectivity, or exchange length scales
therefore require their own $Pe$-resolved assessment.

The resolved-to-DBS comparison tests this practical consequence of the
intrinsic-area representation directly. For this micromodel, the small
inventory differences and modest axial-profile differences show that replacing
explicitly resolved matrix pore walls with cell-supported reactive area does
not materially alter domain-scale Li uptake across the tested $Pe$ range.

Transport and chemistry are coupled by a first-order Lie splitting, which
introduces a splitting error that scales with the timestep and the stiffness of
the reaction terms. For the surface-complexation systems used here, that error
sits within the benchmark agreement reported in Section~\ref{sec:validation},
but it may become significant for faster reactions or stiffer geochemical
systems. A second-order Strang splitting \citep{strang1968construction} is the
natural extension.

Restart depends on the serialised PHREEQC state written with each complete
checkpoint, because the finite-volume \texttt{Surf\_*} fields alone cannot
reconstruct that state. The production $Pe$ sweep restarted from these
checkpoints across several jobs, and moving a run between machines requires the
field data and its matching PHREEQC dump to travel together.

The cases here use a fixed surface area and equilibrium surface complexation,
and cover competitive multi-component chemistry and a Li-selective multi-ion
brine within that scope. Evolving reactive area, mineral dissolution and
precipitation, kinetic surface reactions, and the wider range of mineral
surfaces and aqueous compositions found in natural systems lie outside it.
Extending the model to them will require benchmarks that separate transport
error, surface-area evolution, and reaction-model uncertainty while retaining
cell-level mass-balance checks.
